\documentclass[11pt]{article}
\usepackage[margin=1in]{geometry}
\usepackage{amsmath,amssymb,amsthm}
\usepackage{graphicx}
\usepackage[hidelinks]{hyperref}

\newtheorem{theorem}{Theorem}
\newtheorem{lemma}{Lemma}
\newtheorem{corollary}{Corollary}
\newtheorem{definition}{Definition}
\newtheorem{remark}{Remark}

\newcommand{\dens}{\operatorname{dens}}
\newcommand{\Lip}{\operatorname{Lip}}
\newcommand{\omegaomega}{\omega_\omega}
\newcommand{\R}{\mathbb R}

\title{A Higher-Smoothness Upgrade for the \(L_p\) Branch of the \(c_0(\Gamma)\) Embedding}
\author{agent\_lane\_15}
\date{June 29, 2026}

\begin{document}
\maketitle

\section*{Status}

This packet records a likely valid partial answer to a question in the 2024
paper of Hajek, Johanis, and Schlumprecht on point character and embeddings
into \(c_0(\Gamma)\) \cite{HJS240100831}.  Their Corollary 28 gives, for
subspaces of \(L_p(\mu)\) with \(1<p<\infty\) and density below
\(\omega_\omega\), a bi-Lipschitz embedding
\[
  \Phi:X\longrightarrow c_0(\Gamma)
\]
whose scalar coordinate functions are \(C^1\) smooth.  Immediately after the
corollary, they ask whether for \(p\geq2\) the component functions can have
higher smoothness.

The answer is affirmative for the \(L_p\)-subspace branch in the natural range
allowed by the smoothness of the \(L_p\) norm: \(C^m\) for every integer
\(m<p\), and for every finite \(m\) when \(p\) is an even integer.

\section*{Source Question}

The source passage is on page 11 of arXiv:2401.00831.

\begin{center}
\includegraphics[width=\textwidth]{figures/open_problem_crop.png}
\end{center}

In text form, after proving the \(C^1\)-coordinate embedding corollary, the
authors write:

\begin{quote}
We do not know whether in the corollary above in the case \(p\ge2\) the
component functions can be of higher smoothness.
\end{quote}

\section*{Background: Property \((\ast^m)\)}

\begin{definition}
Let \(m\in\mathbb N\).  A Banach space \(X\) has property \((\ast^m)\) if
there is a constant \(C_0\geq1\) such that for every Lipschitz function
\(f:X\to\R\) and every \(\varepsilon>0\), there is a \(C^m\)-smooth Lipschitz
function \(g:X\to\R\) with
\[
  |f(x)-g(x)|<\varepsilon \quad(x\in X),
  \qquad
  \Lip(g)\leq C_0\Lip(f).
\]
\end{definition}

This is the standard Lipschitz smooth-approximation property used in the
Hajek--Johanis smooth approximation framework.  Fry and Keener prove that a
WCG Banach space with a \(C^m\) smooth norm has property \((\ast^m)\)
\cite[Theorem 4]{FryKeener09070241}.  Their result is stated for \(C^p\)
smoothness; for an integer \(m\) we apply it with \(p=m\).

\section*{Proof Intuition}

The \(c_0(\Gamma)\)-valued part of the source proof is metric: once
\(\dens X<\omega_\omega\), the paper gives a bi-Lipschitz embedding of \(X\)
into some \(c_0(\Gamma)\).  Smoothness of the scalar coordinates is supplied
afterwards by a separate coordinate-smoothing theorem of Hajek and Johanis.

That smoothing theorem does not use a specifically first-order argument.  Its
input is the ability to approximate scalar Lipschitz functions by smooth
Lipschitz functions with controlled Lipschitz constants; the support and
threshold bookkeeping that keeps the map \(c_0\)-valued and bi-Lipschitz is
independent of the differentiability order.  Therefore replacing property
\((\ast^1)\) by \((\ast^m)\) replaces \(C^1\) coordinates by \(C^m\)
coordinates.

For \(L_p\)-subspaces, the remaining input is exactly the classical smoothness
of the \(L_p\) norm together with Fry--Keener's WCG approximation theorem.

\section*{Order-\(m\) Coordinate Smoothing}

\begin{lemma}[Order-\(m\) Hajek--Johanis coordinate smoothing]
\label{lem:coordinate}
Let \(m\in\mathbb N\).  Suppose a Banach space \(X\) has property
\((\ast^m)\).  If \(X\) admits a bi-Lipschitz embedding into some
\(c_0(\Gamma)\), then \(X\) admits a bi-Lipschitz embedding into some
\(c_0(\Gamma')\) whose coordinate functions are \(C^m\) smooth.
\end{lemma}

\begin{proof}
This is the order-\(m\) form of the coordinate-smoothing theorem used by
Hajek--Johanis--Schlumprecht through \cite[Theorem 7.63]{HJBook} in their
proof of the \(C^1\)-coordinate corollary.  We spell out the only point that
depends on differentiability order.

In the Hajek--Johanis construction, a bi-Lipschitz \(c_0(\Gamma)\)-embedding
is first converted into a family of real-valued Lipschitz coordinate tests
with fixed quantitative margins.  These tests are arranged in point-finite
or thresholded families, so that at every \(x\in X\) the resulting coordinate
vector is in \(c_0\), and for every pair \(x,y\) at least one coordinate
separates them by a fixed multiple of \(\|x-y\|\).  The last step replaces
each real-valued Lipschitz test by a smooth Lipschitz approximant whose
uniform error is chosen smaller than the relevant margin.  The upper
Lipschitz bound follows from the common Lipschitz control of the approximants;
the lower Lipschitz bound survives because the approximation errors are below
the separation margins; and the \(c_0\)-valued property survives because the
same point-finite/threshold support bookkeeping is used after the replacement.

For the \(C^1\) statement, this last replacement uses property \((\ast^1)\).
If \(X\) has property \((\ast^m)\), the identical replacement produces
\(C^m\)-smooth approximants with the same Lipschitz and error estimates.
No subsequent step differentiates the approximants or uses a feature special
to \(m=1\).  Thus every coordinate in the final \(c_0(\Gamma')\)-valued map is
\(C^m\), and the same constants give a bi-Lipschitz embedding.
\end{proof}

\section*{Main Partial Result}

\begin{theorem}
\label{thm:lp}
Let \(X\) be a closed subspace of \(L_p(\mu)\), where \(2\leq p<\infty\), and
suppose that \(\dens X<\omega_\omega\).  Let \(m\in\mathbb N\).

If \(m<p\), then \(X\) admits a bi-Lipschitz embedding into some
\(c_0(\Gamma')\) whose coordinate functions are \(C^m\) smooth.  If \(p\) is
an even integer, the same conclusion holds for every finite \(m\).
\end{theorem}

\begin{proof}
By Corollary 26 of \cite{HJS240100831}, \(X\) admits a bi-Lipschitz embedding
into some \(c_0(\Gamma)\), because \(X\subset L_p(\mu)\) and
\(\dens X<\omega_\omega\).

Since \(1<p<\infty\), the space \(L_p(\mu)\) is reflexive, and every closed
subspace \(X\) is reflexive.  In particular \(X\) is weakly compactly
generated.  For every integer \(m<p\), the usual \(L_p\)-norm is \(C^m\)
smooth away from the origin; when \(p\) is an even integer, the corresponding
power function is polynomial and the usual smooth-norm argument gives
\(C^m\) smoothness for every finite \(m\).  Restricting the norm to \(X\)
therefore gives \(X\) a \(C^m\) smooth norm in the stated range.

Fry--Keener \cite[Theorem 4]{FryKeener09070241} now gives property
\((\ast^m)\) for \(X\).  Applying Lemma \ref{lem:coordinate} to the
bi-Lipschitz \(c_0(\Gamma)\)-embedding supplied by the source paper yields a
bi-Lipschitz embedding of \(X\) into another \(c_0\)-space with \(C^m\) smooth
coordinate functions.
\end{proof}

\begin{corollary}
For every closed subspace \(X\subset L_p(\mu)\) with \(2\leq p<\infty\) and
\(\dens X<\omega_\omega\), the embedding in the source corollary can be chosen
with coordinate functions smoother than \(C^1\).  More precisely, one can take
all coordinates \(C^2\) when \(p>2\), all finite orders when \(p=2\), and all
orders \(m<p\) in general.
\end{corollary}

\begin{proof}
If \(p=2\), apply Theorem \ref{thm:lp} in the even-integer case.  If \(p>2\),
take \(m=2\), or any larger integer \(m<p\) when available.
\end{proof}

\section*{Verification Notes}

The proof uses three external inputs.

First, the source paper supplies the underlying bi-Lipschitz embedding of
\(L_p\)-subspaces of density below \(\omega_\omega\) into \(c_0(\Gamma)\)
\cite[Corollary 26]{HJS240100831}.  Second, Hajek--Johanis coordinate
smoothing is used in the same way the source uses it for the \(C^1\) corollary;
the only change is replacing property \((\ast^1)\) by property \((\ast^m)\).
Third, Fry--Keener supplies property \((\ast^m)\) for WCG spaces with a
\(C^m\) smooth norm \cite[Theorem 4]{FryKeener09070241}.

No numerical or symbolic computation is involved.

\section*{Limitations}

This is a partial answer to the source question.  It covers the \(L_p\) branch
of the source corollary in the differentiability range supported by the
\(L_p\) norm.  It does not settle the super-reflexive Banach-lattice branch
unless the lattice in question is separately known to satisfy property
\((\ast^m)\).  It also does not claim \(C^m\) coordinates for non-even
\(L_p\)-spaces in orders \(m\geq p\).

\section*{Novelty Check}

The run indexes were searched for arXiv:2401.00831, smooth-coordinate
keywords, \(C^m\)/higher-smoothness phrases, Fry--Keener, point character, and
\(c_0(\Gamma)\) embedding terminology.  The only prior run item was the
conditional packet isolating this same route.  A bounded web search for the
source title together with ``higher smoothness'', ``component functions'', and
``Theorem 7.63'' did not reveal a later explicit answer.  The present packet
therefore records the route as new run progress, subject to human review of
the order-\(m\) reading of the Hajek--Johanis coordinate-smoothing theorem.

\section*{Human Review Recommendation}

Review Lemma \ref{lem:coordinate} first.  The substantive check is that the
Hajek--Johanis coordinate-smoothing theorem used by the source is indeed
order-free once property \((\ast^m)\) is available.  After that, the
Fry--Keener and \(L_p\)-smoothness inputs are standard.

\begin{thebibliography}{9}

\bibitem{HJS240100831}
P. Hajek, M. Johanis, and T. Schlumprecht,
\emph{Remarks on the point character of Banach spaces and non-linear
embeddings into \(c_0(\Gamma)\)},
arXiv:2401.00831, 2024.

\bibitem{HJBook}
P. Hajek and M. Johanis,
\emph{Smooth analysis in Banach spaces},
De Gruyter Series in Nonlinear Analysis and Applications, vol. 19,
De Gruyter, Berlin, 2014.

\bibitem{FryKeener09070241}
R. Fry and L. Keener,
\emph{Approximation of Lipschitz functions by Lipschitz, \(C^p\) smooth
functions on weakly compactly generated Banach spaces},
arXiv:0907.0241, 2009.

\end{thebibliography}

\end{document}
